\documentclass[UTF8,a4paper,zihao=-4]{ctexart}

\usepackage[left=2.5cm,right=2.5cm,top=2.6cm,bottom=2.6cm]{geometry}
\usepackage{amsmath,amssymb}
\usepackage{booktabs}
\usepackage{array}
\usepackage{multirow}
\usepackage{float}
\usepackage{siunitx}
\usepackage{enumitem}
\usepackage{setspace}
\usepackage{fontspec}

\setmainfont{Times New Roman}
\setCJKmainfont{FandolSong-Regular}
\sisetup{detect-all}

\setstretch{1.5}
\setlength{\parindent}{2em}
\pagestyle{plain}

\AtBeginDocument{\songti}

\begin{document}

\begin{titlepage}
    \centering
    {\zihao{2}\bfseries 实验报告\par}
    \vspace{1cm}
    {\zihao{-2}\bfseries 实验5.21\quad 金属逸出功实验\par}
    \vspace{2cm}

    \renewcommand{\arraystretch}{1.8}
    \begin{tabular}{p{4cm}p{8cm}}
        课程名称： & 大学物理实验 \\
        姓名： & 【待补充】 \\
        学号： & 【待补充】 \\
        专业班级： & 【待补充】 \\
        开课学期： & 【待补充】 \\
        实验日期： & 【待补充】 \\
    \end{tabular}

    \vfill
    {\zihao{4} 物理与光电学院公共教学实验中心\par}
    \vspace{0.5cm}
    {\zihao{4} 【待补充】年【待补充】月\par}
\end{titlepage}

\newpage

\section*{一、实验目的}

用 Richardson 直线法测定金属钨的电子逸出功。

\section*{二、实验原理}

金属中的电子在热运动过程中，如果其动能足以克服表面势垒，就能逸出金属表面形成热电子发射。根据 Richardson-Dushman 方程，热电子发射电流密度为：

\begin{equation}
    J = A T^2 \exp\left(-\frac{eW}{kT}\right)
\end{equation}

其中 $J$ 为发射电流密度，$A$ 为 Richardson 常数，$T$ 为绝对温度，$W$ 为电子逸出功，$e$ 为元电荷，$k$ 为 Boltzmann 常数。

对于实际阴极，发射电流为：

\begin{equation}
    I = A S T^2 \exp\left(-\frac{eW}{kT}\right)
\end{equation}

其中 $S$ 为阴极发射面积。

取对数得：

\begin{equation}
    \lg\left(\frac{I}{T^2}\right) = \lg(AS) - \frac{eW}{k\ln 10} \cdot \frac{1}{T}
\end{equation}

以 $\lg(I/T^2)$ 为纵坐标，$1/T$ 为横坐标作图，可得一条直线，其斜率为：

\begin{equation}
    m = -\frac{eW}{k\ln 10}
\end{equation}

由此可求得逸出功：

\begin{equation}
    W = -m \cdot \frac{k}{e} \cdot \ln 10
\end{equation}

由于实际测量中存在接触电势差和空间电荷效应，需要在不同阳极电压 $U_a$ 下测量阳极电流 $I_a$，然后外推到 $U_a = 0$ 得到零场发射电流 $I$。

根据 Schottky 效应，$\lg I_a$ 与 $\sqrt{U_a}$ 成线性关系：

\begin{equation}
    \lg I_a = a\sqrt{U_a} + \lg I
\end{equation}

通过线性拟合 $\lg I_a$ - $\sqrt{U_a}$ 关系，可得截距 $\lg I$，即零场电流的对数。

\section*{三、实验仪器}

理想二极管实验装置、直流稳压电源、数字万用表等。

\section*{四、实验内容与数据记录}

\subsection*{1. 测量不同灯丝电流下的阳极电流}

在灯丝电流 $I_f = 0.600\sim0.700\,\text{A}$ 范围内，以 $0.025\,\text{A}$ 为间隔调节灯丝电流，对每个 $I_f$ 值，测量不同阳极电压 $U_a$ 下的阳极电流 $I_a$（单位：$\mu\text{A}$）。

温度计算公式：$T = 900 + 1430 I_f$（单位：K）

\subsection*{2. 原始数据记录}

\begin{table}[H]
    \centering
    \caption{附录表5.21.1 不同 $U_a$ 和 $I_f$ 下的 $I_a$ 数据（单位：$\mu\text{A}$）}
    \renewcommand{\arraystretch}{1.3}
    \begin{tabular}{c|ccccccc}
        \hline
        $U_a/\text{V}$ & 0.16 & 0.25 & 0.36 & 0.49 & 0.64 & 0.81 & 1.00 \\
        \hline
        $I_f = 0.600\,\text{A}$ & 59 & 62 & 66 & 69 & 68 & 72 & 76 \\
        $I_f = 0.625\,\text{A}$ & 110 & 117 & 125 & 127 & 131 & 137 & 142 \\
        $I_f = 0.650\,\text{A}$ & 199 & 211 & 221 & 230 & 238 & 244 & 252 \\
        $I_f = 0.675\,\text{A}$ & 336 & 365 & 383 & 399 & 410 & 423 & 437 \\
        $I_f = 0.700\,\text{A}$ & 544 & 603 & 637 & 670 & 693 & 719 & 736 \\
        \hline
    \end{tabular}
\end{table}

\section*{五、数据处理}

\subsection*{1. 计算各灯丝电流对应的温度}

\begin{table}[H]
    \centering
    \caption{灯丝电流与温度对应关系}
    \renewcommand{\arraystretch}{1.3}
    \begin{tabular}{ccccc}
        \hline
        $I_f/\text{A}$ & 0.600 & 0.625 & 0.650 & 0.675 & 0.700 \\
        \hline
        $T/\text{K}$ & 1758.0 & 1793.8 & 1829.5 & 1865.2 & 1901.0 \\
        \hline
    \end{tabular}
\end{table}

\subsection*{2. $\lg I_a$ - $\sqrt{U_a}$ 线性拟合}

对每个温度下（即每个 $I_f$ 值）的 $\lg I_a$ 与 $\sqrt{U_a}$ 进行线性拟合，得到截距 $\lg I$（零场电流对数）。

\begin{table}[H]
    \centering
    \caption{$\lg I_a$ - $\sqrt{U_a}$ 线性拟合结果}
    \renewcommand{\arraystretch}{1.3}
    \begin{tabular}{cccccc}
        \hline
        $I_f/\text{A}$ & $T/\text{K}$ & 斜率 & $\lg I$ & 相关系数 $r$ \\
        \hline
        0.600 & 1758.0 & 0.1688 & 1.7093 & 0.9753 \\
        0.625 & 1793.8 & 0.1750 & 1.9798 & 0.9867 \\
        0.650 & 1829.5 & 0.1664 & 2.2399 & 0.9904 \\
        0.675 & 1865.2 & 0.1786 & 2.4682 & 0.9799 \\
        0.700 & 1901.0 & 0.2083 & 2.6700 & 0.9728 \\
        \hline
    \end{tabular}
\end{table}

\subsection*{3. Richardson 直线法求逸出功}

根据 Richardson-Dushman 方程，以 $\lg(I/T^2)$ 为纵坐标，$1/T$ 为横坐标作直线拟合。

\begin{table}[H]
    \centering
    \caption{Richardson 直线法数据}
    \renewcommand{\arraystretch}{1.3}
    \begin{tabular}{ccccc}
        \hline
        $T/\text{K}$ & $\lg I$ & $\lg(I/T^2)$ & $10^4/T\,\text{K}^{-1}$ \\
        \hline
        1758.0 & 1.7093 & $-4.7808$ & 5.688 \\
        1793.8 & 1.9798 & $-4.5277$ & 5.575 \\
        1829.5 & 2.2399 & $-4.2848$ & 5.466 \\
        1865.2 & 2.4682 & $-4.0733$ & 5.361 \\
        1901.0 & 2.6700 & $-3.8880$ & 5.260 \\
        \hline
    \end{tabular}
\end{table}

线性拟合结果：

\begin{equation}
    \lg\left(\frac{I}{T^2}\right) = -20962.8553 \times \frac{1}{T} + 7.1561
\end{equation}

相关系数 $r = -0.9992$

\subsection*{4. 逸出功计算}

斜率 $m = -20962.8553$

\begin{equation}
    W = -m \cdot \frac{k}{e} \cdot \ln 10 = 20962.8553 \times 8.6173 \times 10^{-5} \times 2.3026 = 4.1595\,\text{eV}
\end{equation}

\section*{六、实验结果}

金属钨的电子逸出功测量值为：

\begin{equation}
    \boxed{W = 4.1595\,\text{eV}}
\end{equation}

钨的逸出功理论值约为 $4.5\,\text{eV}$，测量值与理论值基本吻合。

\end{document}
